\section*{Lernziele: Zahlensysteme und Kodierungen}
\begin{todolist}
\item Ich kann Zahlen von einem gegebenen Zahlensystem in das Dezimalsystem umwandeln. Beispiele
\begin{itemize}
\item $101_2 = 5_{10}$
\item $204_5 = 54_{10}$
\item $1A_{16} = 26_{10}$
\end{itemize}
\item Ich kann Zahlen vom Dezimalsystem in ein vorgegebenes Zahlensystem umwandeln. Beispiele:
\begin{itemize}
\item $5_{10} = 101_2$
\item $135_{10}=1020_5$
\item $26_{10}=1A_{16}$
\end{itemize}
\item Ich kann Zahlen in einem fremden Zahlensystem (schriftlich) addieren. Zum Beispiel:
\[123_5 + 431_5 = 1104_5\]
\item Ich kenne die Einheiten Bits und Bytes und kann Beispiele nennen, wie diese physikalisch realisiert werden können.
\item Ich kann berechnen, wie viele unterschiedlichen Informationen mit einer gewissen Anzahl Bits codiert werden können.
\item Ich kann erklären, worum es sich bei ASCII handelt (Sie müssen die Tabelle aber nicht auswendig kennen \emoji{winking-face}).
\item Ich verstehe Bilder als ein Raster aus Pixeln, die je eine bestimmte Farbe haben.
\item Ich kenne den Begriff der Farbtiefe und kann berechnen, wie viele Farben pro Pixel bei einer gegebenen Farbtiefe möglich sind.
\item Ich kann berechnen, wie viele Bits nötig sind, um ein Bild mit einer bestimmten Anzahl möglichen Farben zu speichern.
\item Ich kann erklären, wie man eine RGB-Farbcodierung wie \texttt{\#1243A3} interpretieren kann.
\item Ich kann einen hexadezimalen Farbcode wie \texttt{\#1243A3} einer Farbe (aus einer Auswahl) zuordnen.

\end{todolist}